%% kapa: exam=1 round=28 year=2017 period=1 form=A subject=real_estate subject_ko=부동산학원론 seq=9 number=89
%% source: Q-Net 감정평가사 기출문제 (cst00309) artlSeq=5199603 file=extracted/1008/5199603/2017년 제28회 감정평가사 1차시험 문제지/1교시 A형.pdf
%% conversion: kordoc 4.24.51 -> tools/gen1.py (auto); review=proofread-sonnet; flags=none
\begin{question}{89}
부동산투자 의사결정방법에 관한 설명으로 옳지 않은 것은?
\begin{choices}
  \item 수익성지수법은 투자된 현금유출의 현재가치를 이 투자로부터 발생되는 현금유입의 현재가치로 나눈 것이다.
  \item 회계적이익률법에서는 상호배타적인 투자안일 경우에 목표이익률보다 큰 투자안 중에서 회계적 이익률이 가장 큰 투자안을 선택한다.
  \item 순현가법은 화폐의 시간가치를 고려한 방법으로 순현가가 “0”보다 작으면 그 투자안을 기각한다.
  \item 내부수익률은 투자안의 순현가를 “0”으로 만드는 할인율을 의미하며, 투자자 입장에서는 최소한의 요구수익률이기도 하다.
  \item 회수기간법은 화폐의 시간적 가치를 고려하지 않고, 회수기간이 더 짧은 투자안을 선택하는 투자결정법이다.
\end{choices}
\end{question}
\answer{1}
